\ifx\allfiles\undefined
\documentclass[12pt, a4paper, oneside, UTF8]{ctexbook}
\def\path{../config}
\input{\path/package.tex}

% 定理环境设置
\input{\path/theorem.tex}

\input{\path/custom.tex}

\input{\path/cover_package} 

% 封面页设置
\def\myTitle{量子力学：理论与应用
}
\def\mySubtitle{Quantum Mechanics: Theory and Application}
\def\myAuthor{Max}
\def\myCreateDateCover{2026年10月11日}
\def\myUpdateDateCover{ \today}
\def\myDateForeword{}
% 前言页设置
\def\myForeword{前言}
\def\myForewordText{

}
\def\mySubheading{}
\begin{document}
% \input{\path/cover_text.tex}

%%%%%%%% 前言页 %%%%%%%%

\newpage
\thispagestyle{empty}
\begin{center}
    \Huge\textbf{\myForeword}
\end{center}
\myForewordText
\begin{flushright}
    \begin{tabular}{c}
        \myDateForeword
    \end{tabular}
\end{flushright}

% 目录页
\newpage
\pagestyle{plain}
\setcounter{page}{1}
\pagenumbering{Roman}
\tableofcontents

% 正文页
\newpage
\pagenumbering{arabic}
\setcounter{chapter}{-1}
\setcounter{page}{1}

\pagestyle{fancy}
% \fancyhead[L]{量子力学笔记}
% \fancyhead[C]{中间页眉}
% \fancyhead[R]{右侧页眉}
% \fancyfoot[C]{\thepage}
\fancyhf{}
% 设置页眉
\fancyhead[L]{\leftmark}    % 左侧：文章标题（\leftmark 通常是章标题）
\fancyhead[C]{ }                       % 中间：留空（或自定义内容）
\fancyhead[R]{\rightmark}   % 右侧：节标题（\rightmark 通常是节标题）


% 设置页脚（示例：中间显示页码）
\fancyfoot[C]{\thepage}

\renewcommand{\headrulewidth}{0.4pt}
\renewcommand{\footrulewidth}{0pt}








\else
\fi

\chapter{一维定态问题}
\begin{examct}{}{}
    一维定态的性质、方势阱、一维谐振子、一维散射问题
\end{examct}
\begin{examrqm}{}{}
    \begin{enumerate}
        \item 掌握一维定态的基本性质、方势阱与一维谐振子的本征值与本征态，了解一维散射问题
        \item 会灵活运用所学知识解决一维定态问题
    \end{enumerate}
\end{examrqm}

\section{一维定态的性质}
\section{方势阱}
\subsection{一维无限深方势阱}
势 $\displaystyle V(x)=\begin{cases}0,&0\leqslant x \leqslant a \\ \infty ,&x<0, x>a \end{cases}$ 
\hfill
\adjustbox{valign=m}{
    \begin{tikzpicture}[>=stealth, % 
                    line width=0.7pt] % 设置线条宽度
    \draw[->] (-2,0) -- (4,0) node[right] {$x$};
    \draw[->] (0,-0.5) -- (0,4) node[above] {$V(x)$};

    \node at (-0.3,-0.3) {$O$};

    \draw (0,3) -- (-1.5,3);
    \node at (-1,3.3) {$\infty$};

    \draw (0,0) -- (2,0); % 势阱内部 (0 < x < a)，V=0
    \draw (2,0) -- (2,3) -- (4,3); % 右侧势垒 (x > a)，V=∞
    \node at (2,-0.3) {$a$}; 
    \node at (3,3.3) {$\infty$};

    \draw (0,0) -- (0,3);

    \node[below] at (1,0) {$V=0$};
\end{tikzpicture}
}

当 $x<0, x>a$ 时, $V(x)=\infty$,$\psi(x)=0$.

当 $0\leqslant x \leqslant a$ 时,$V(x)=0$,有
\begin{equation}
    -\frac{\hbar^2}{2m}\frac{d^2}{dx^2}\psi(x)=E\psi(x)
\end{equation}
波函数为
\begin{equation}
    \psi_n(x)=\sqrt{\frac{2}{a}}\sin(\frac{n\pi x}{a}),\quad n=1,2,3,\cdots
\end{equation}
能量为
\begin{equation}
    E_n=\frac{n^2\pi^2\hbar^2}{2ma^2},\quad n=1,2,3,\cdots
\end{equation}

详细推导见\ref{chap:1DInfiniteSquarePotentialWell}

一维无限深方势阱性质：
\begin{enumerate}
    \item 能量本征函数时分立的，且正交归一。$\displaystyle \braket{\psi_n}{\psi_m}=$
\end{enumerate}

\subsection{三维有限深方势阱}
势 $\displaystyle V(x)=\begin{cases}
    0,&0<x<a,0<y<a,0<z<a\\
    \infty,&other
\end{cases}$



\subsection{一维有限深方势阱}
势 $\displaystyle V(x)=\begin{cases}
    0,&-a<x<a\\
    V_0,&|x|>a
\end{cases}$
\hfill
\adjustbox{valign=m}{  % valign=m 使图片中线与基线对齐
    \begin{tikzpicture}
        \draw[->] (-3,0) -- (3,0) node[right] {$x$};
        \draw[->] (0,-0.5) -- (0,4) node[above] {$V(x)$};
        \node at (-0.3,-0.3) {$O$};
        \draw (-1,0) -- (-1,3) -- (-3,3);
        \node at (-1,-0.3) {$-a$};
        \node at (-3,3.3) {$V_0$};

        \draw (1,0) -- (1,3) -- (3,3);
        \node at (1,-0.3) {$a$};
        \node at (3,3.3) {$V_0$};
    \end{tikzpicture}
}

\begin{tikzpicture}[>=stealth,scale=1.5]
    \draw[->] (-2,0) -- (2,0) node[right] {$x$};
    \draw[->] (0,-0.3) -- (0,4) node[above] {$V(x)$};

    \node at (-0.3,-0.3) {$O$};

    \draw[thick, domain=-3:3, samples=300] plot (\x, {exp(-(\x)^2/(2*0.11^2))/(0.11*sqrt(2*pi))});

    \node at (0.5, 3) {$\delta(x)$};

\end{tikzpicture}


\section{一维谐振子}
\subsection{谐振子势}
$\displaystyle V(x)=\frac{1}{2}m\omega^2x^2$
\hfill
\adjustbox{valign=m}{
    \begin{tikzpicture}[>=stealth, % 
                    line width=0.7pt] % 设置线条宽度
    \draw[->] (-2.5,0) -- (2.5,0) node[right] {$x$};
    \draw[->] (0,-0.5) -- (0,3) node[above] {$V(x)$};
    \node at (-0.3,-0.3) {$O$};

    \draw[thick, domain=-2:2, samples=300] plot (\x, {0.5*(\x)^2});

\end{tikzpicture}
}

能量为：
\begin{equation}
    E_n=(n+\frac{1}{2})\hbar\omega \quad n=1,2,3,\cdots
\end{equation}

波函数为
\begin{equation}
    \psi_n(x)=N_ne^{-\alpha^2x^2/2}H_n(\alpha x)
\end{equation}
其中$\displaystyle \alpha=\sqrt{\frac{m\omega}{\hbar}}$\ ,\ $\displaystyle N_n=\sqrt{\frac{\alpha}{\pi^{\frac{1}{2}}2^n n!}}$为归一化系数，$H_n(\alpha x)$为 Hermite 多项式。

基态波函数为
\begin{equation}
    \psi_0(x)=N_0e^{-\alpha^2x^2/2}
\end{equation}

详细推导见\ref{chap:1DHarmonicOscillator}







\section{一维散射问题}
\section{附录}
\subsection{一维无限深方势阱推导 \label{chap:1DInfiniteSquarePotentialWell}}
一维无限深方势阱势函数为
\begin{equation}
    V(x)=\begin{cases}
            0,&0\leqslant x \leqslant a \\ 
            \infty ,&x<0, x>a 
        \end{cases}
\end{equation} 

当 $x<0, x>a$ 时, $V(x)=\infty$,$\psi(x)=0$.

当 $0\leqslant x \leqslant a$ 时,$V(x)=0$,有
\begin{equation}
    -\frac{\hbar^2}{2m}\frac{d^2}{dx^2}\psi(x)=E\psi(x)
\end{equation}
令 $\displaystyle k=\sqrt{\frac{2mE}{\hbar^2}}$ 整理得 
\begin{equation}
    \frac{d^2}{dx^2}\psi(x) = -\frac{2mE}{\hbar^2}\psi(x) = -k^2\psi(x)
\end{equation}
解得
\begin{equation}
    \psi(x)=Ae^{ikx}+Be^{-ikx}
\end{equation}
又该势为驻波解，可写为三角形式
\begin{equation}
    \psi(x)=A\sin(kx)+B\cos(kx)
\end{equation}
代入边界条件$\begin{cases} \psi(0)=0\\ \psi(a)=0 \end{cases}$ 得
\begin{equation}
    \begin{cases}
        B = 0 \\
        A\sin(ka)=0
    \end{cases}
\end{equation}
对于 $A\sin(ka)=0$ ,当 $A=0$ 时，$\psi(x)=0$，无法归一化。所以
\begin{equation}
    \sin(ka)=0 \quad \Rightarrow \quad k=\frac{n\pi}{a} \quad n=1,2,3,\cdots
\end{equation}
所以
\begin{equation}
    \psi_n(x) = A \sin\left(\frac{n\pi x}{a}\right)
\end{equation}
又 $\displaystyle k=\frac{n\pi}{a} = \sqrt{\frac{2mE}{\hbar^2}}$ 得到
\begin{equation}
    \boxed{E_n=\frac{n^2\pi^2\hbar^2}{2ma^2} \quad n=1,2,3,\cdots}
\end{equation}
归一化
\begin{equation}
    \begin{aligned}
        1 &=\int_{-\infty}^{+\infty}\left| \psi_n(x) \right|^2 dx  \\
        &=\int_{0}^a \psi_n^*(x)\psi_n(x) dx  \\
        &= \left|A\right|^2 \int_{0}^a \sin^2\left(\frac{n\pi x}{a}\right) dx\\
        &= \left|A\right|^2 \int_{0}^a \frac{1}{2}\left(1-\cos(\frac{2n\pi x}{a})\right) dx\\
        &= \left|A\right|^2 \left( \frac{a}{2} - \int_{0}^a\cos(\frac{2n\pi x}{a}) dx\right)\\
        &= \frac{a}{2}\left|A\right|^2
    \end{aligned}
\end{equation}
得 
\begin{equation}
    A = \sqrt{\frac{2}{a}}
\end{equation}
所以定态波函数为
\begin{equation}
    \boxed{\Psi(x,t) = \sqrt{\frac{2}{a}} \sin\left(\frac{n\pi x}{a}\right) e^{-iE_nt/\hbar}}
\end{equation}


\subsection{一维谐振子势推导 \label{chap:1DHarmonicOscillator}}
一维谐振子势为：
\begin{equation}
    V(x)=\frac{1}{2}m\omega^2x^2
\end{equation}

由于 $ V_{\pm\infty} = \infty $ , $ E<V_{\pm\infty}$ ,所以只存在束缚态。




\ifx\allfiles\undefined
\end{document}
\fi